\chapter{Conclusions and Future Work}

\label{chap:Chapter7}

\section{Conclusions}
\label{sec:conclusions}

This dissertation investigated an incremental approach for migrating selected low-level Boost functionality from C++ to Rust while preserving interoperability with existing C++ software. The work treated migration as an architectural and semantic transformation rather than as a direct source translation. Two case studies, \texttt{Boost.Align} and \texttt{Boost.Endian}, were implemented using different interoperability strategies so that the methodology could be evaluated under different compatibility constraints.

The study demonstrates that selected functionality from both libraries can be expressed using Rust-native abstractions while retaining a usable C++ integration boundary. The implementations also show that the appropriate FFI mechanism depends on the compatibility requirements of the component: the Align implementation uses a small C-compatible interface because the target use case retains C++03 compatibility, whereas the Endian implementation uses the \texttt{cxx} bridge.

The evaluation provides quantitative evidence for the selected workloads. In the Boost.Align end-to-end benchmark, the Rust implementation measured 20.7014~ms per iteration compared with 21.6782~ms for the Boost implementation, corresponding to approximately 4.7\% lower measured execution time. The same experiment reported a larger RSS increase during allocation for the Rust implementation, so the results do not support a general claim of reduced memory consumption. In the Boost.Endian benchmark, Rust through the \texttt{cxx} bridge measured lower execution time than the manual C++ byte-manipulation baseline at all four tested input sizes, with a reported average speedup of approximately 35\%. Because that baseline is not the original Boost.Endian implementation, the result should not be interpreted as evidence that Rust is generally faster than Boost.Endian.

The central conclusion is therefore methodological: successful incremental migration depends on identifying observable semantics, selecting an appropriate Rust abstraction, defining an explicit interoperability contract, isolating unsafe operations, and validating each migrated component with evidence appropriate to the property being claimed. The two case studies demonstrate these patterns without establishing that every Boost component, particularly highly template-dependent components, can be migrated using the same implementation choices.

\section{Answers to the Research Questions}
\label{sec:answers_research_questions}

\subsection{RQ1 -- Architectural and interoperability patterns}

The implementation identifies several reusable patterns: semantic rather than syntactic translation, explicit ownership and lifetime contracts, small FFI surfaces, isolation of unsafe operations, compatibility-driven FFI selection, and validation against observable behaviour. These patterns are implemented in different forms for Align and Endian rather than through a single universal binding mechanism.

\subsection{RQ2 -- Component-selection criteria}

The methodology defines selection criteria covering automotive relevance, safety relevance, isolation capability, interoperability constraints, portability, performance relevance, implementation complexity, and the potential to expose reusable migration patterns. These criteria provide a traceable basis for selecting case studies rather than selecting components solely because they are easy to translate.

\subsection{RQ3 -- Rust reimplementation while preserving behaviour}

The two case studies demonstrate that selected Boost functionality can be re-expressed using Rust-native abstractions while preserving the tested observable behaviour. The implementations do not reproduce the complete Boost APIs or their internal C++ structures. Consequently, the evidence supports semantic reproduction of the selected functionality rather than complete API equivalence.

\subsection{RQ4 -- FFI boundary design}

The implementations show two compatible approaches. Boost.Align uses a small C-compatible boundary suitable for C++03 consumers, while Boost.Endian uses \texttt{cxx}. In both cases, interface design requires explicit reasoning about representation, ownership, lifetime, valid inputs, and error or panic behaviour. The appropriate mechanism is therefore determined by the compatibility requirements of the component rather than selected independently of them.

\subsection{RQ5 -- Observed performance characteristics}

The evaluated workloads show competitive measured execution times for the Rust implementations. Boost.Align measured approximately 4.7\% lower execution time in the selected end-to-end AVX workload, while Boost.Endian measured approximately 35\% lower execution time than the manual C++ baseline across the tested sizes. These are workload-specific measurements and do not establish universal performance improvements. In particular, the Align implementations did not use identical implementation details in every aspect, and the Endian baseline was not the original Boost.Endian implementation.

\subsection{RQ6 -- Safety guarantees and FFI contracts}

Within safe Rust, ownership, borrowing, bounds, and representation invariants can be expressed and checked by the language and type system. At the FFI boundary, these guarantees become dependent on interface contracts and foreign-code behaviour. The caller must respect requirements concerning pointer validity, buffer length, alignment, lifetime, aliasing, and data layout. The migration therefore reduces the manually managed safety surface inside the Rust component but does not remove the need for verification of the hybrid boundary.

\section{Contributions}
\label{sec:contributions}

The dissertation contributes a structured component-selection framework, a migration methodology, two complementary low-level case studies, an analysis of generalizable versus library-specific migration decisions, and an experimental evaluation of correctness, performance, memory observations, and FFI overhead. It also makes explicit the distinction between properties enforced by Rust internally and properties that remain contracts across FFI.

\section{Limitations}
\label{sec:limitations}

The work is limited to selected functionality from two Boost libraries and therefore does not establish complete migration of either library or the feasibility of migrating the entire Boost ecosystem. The experiments were conducted on a limited set of workloads and hardware configurations. The performance measurements characterize observed execution time and throughput but do not establish worst-case execution time, jitter, schedulability, or hard real-time guarantees. The memory measurements are process-level RSS observations rather than a complete allocator or memory-efficiency analysis.

The safety discussion likewise should not be interpreted as a certification argument. The dissertation identifies where Rust's static guarantees apply and where explicit FFI contracts remain necessary, but it does not provide a complete ISO~26262 safety case, formal proof, or certification assessment of the resulting software.

Finally, the Endian performance comparison uses a manual C++ byte-manipulation baseline rather than a direct Boost.Endian implementation. The Align benchmark also contains implementation differences that limit the interpretation of the measured performance difference. These limitations are considered explicitly when drawing conclusions from the results.

\section{Future Work}
\label{sec:future_work}

The following sections consolidate the implementation improvements and validation activities identified during the study. They are future work rather than evidence required to support the conclusions above.

\section{Boost.Align Improvements}

\label{sec:align-improvements}

The Boost.Align implementation described in Chapter~\ref{chap:Chapter6} provides an aligned
buffer abstraction, compile-time alignment through Rust's
\texttt{\#[repr(align(...))]} attribute, AVX-based processing, and a
C-compatible interface. The implementation successfully demonstrates that
the alignment requirements of the selected use case can be expressed using
Rust's type system.

The current implementation is nevertheless intentionally narrower than the
complete Boost.Align API. The following improvements would increase its
generality and suitability for production use.


% ------------------------------------------------------------
\subsection{Generalisation of Alignment Requirements}

The current \texttt{Aligned<T>} abstraction uses a fixed 32-byte alignment:

\begin{verbatim}
#[repr(align(32))]
pub struct Aligned<T>(pub T);
\end{verbatim}

This is appropriate for the AVX workload evaluated in this project because
a 256-bit AVX register contains eight 32-bit floating-point values and the
aligned load and store operations used by the implementation require
32-byte-aligned memory.

However, fixing the alignment to 32 bytes limits the abstraction to one
particular use case. A more general implementation should allow the required
alignment to be selected according to the type or operation being performed.

The existing \texttt{aligned\_type!} macro partially addresses this
limitation by allowing a user-defined alignment:

\begin{verbatim}
aligned_type!(TypeName, InnerType, Alignment);
\end{verbatim}

A future implementation could replace the macro-based interface with a more
general abstraction based on Rust's type system and allocation facilities.
Such an abstraction would allow the same implementation to support different
alignment requirements where required.

This would make the Rust implementation more closely applicable to the wider
class of alignment requirements supported by Boost.Align, whose purpose is to
provide portable control and inspection of memory alignment
\cite{boostalign}.


% ------------------------------------------------------------
\subsection{Expansion of the Boost.Align API}

The current implementation focuses on aligned storage and aligned SIMD
processing rather than reproducing the complete Boost.Align API.

In particular, the implementation does not currently provide direct Rust
equivalents for all of the following Boost.Align facilities:

\begin{itemize}
    \item alignment calculation;
    \item \texttt{align\_up};
    \item \texttt{align\_down};
    \item alignment inspection;
    \item \texttt{is\_aligned};
    \item general aligned allocation and deallocation;
    \item aligned allocators; and
    \item allocator adaptors.
\end{itemize}

These facilities could be added as separate Rust abstractions rather than as
a direct translation of the original C++ API.

For example, alignment calculations could be expressed using integer
operations, while allocation could be represented using Rust's
\texttt{Layout} abstraction. The objective should remain to preserve the
semantic guarantees of the original functionality rather than reproduce its
C++ interface literally.

This would also allow a more direct comparison between the facilities
provided by Boost.Align and their Rust equivalents.


% ------------------------------------------------------------
\subsection{Removal of Unnecessary Raw Slice Construction}

The current \texttt{AlignedBuffer} implementation provides the following
operations:

\begin{verbatim}
pub fn as_slice(&self) -> &[f32]
pub fn as_mut_slice(&mut self) -> &mut [f32]
\end{verbatim}

These operations construct Rust slices from raw pointers using
\texttt{std::slice::from\_raw\_parts} and
\texttt{std::slice::from\_raw\_parts\_mut}. Consequently, they require
\texttt{unsafe} code even though the underlying buffer is owned by a safe
Rust abstraction.

A future implementation should investigate whether the internal
representation can be reorganised so that the safe slice interface does not
require raw-pointer reconstruction.

For example, the internal representation could use an aligned allocation
whose ownership and length are represented explicitly while exposing a safe
slice through a carefully designed abstraction. The goal would be to reduce
the amount of unsafe code while preserving the alignment guarantee.

This would strengthen one of the central objectives of the migration
strategy: minimizing and isolating unsafe Rust code.


% ------------------------------------------------------------
\subsection{Improved FFI Error Handling}

The current FFI implementation uses assertions for invalid pointers and
invalid buffer lengths. For example, the FFI functions verify that a buffer
pointer is not null and that the buffer length is compatible with the AVX
operation.

Assertions are useful during development because they expose violations of
the expected preconditions immediately. However, they are not an ideal
error-handling mechanism for a production C++ interface.

A future interface should replace assertions at the FFI boundary with
explicit error handling. Possible approaches include returning an integer
status code, returning a nullable pointer where appropriate, or exposing a
dedicated error enumeration.

For example:

\begin{verbatim}
enum AlignError {
    NullPointer,
    InvalidLength,
    UnsupportedArchitecture,
    InvalidAlignment
}
\end{verbatim}

Such an interface would allow the C++ caller to handle invalid input without
depending on Rust panic behavior.

This is particularly important for safety-oriented systems, where error
handling should be deterministic and explicitly defined at component
boundaries.


% ------------------------------------------------------------
\subsection{Runtime SIMD Feature Detection}

The current SIMD implementation is explicitly compiled with:

\begin{verbatim}
#[target_feature(enable = "avx")]
\end{verbatim}

and uses AVX intrinsics such as
\texttt{\_mm256\_load\_ps},
\texttt{\_mm256\_add\_ps}, and
\texttt{\_mm256\_store\_ps}.

This is suitable for the experimental environment used by the project but
introduces an architectural assumption: the target processor must support
AVX.

A more portable implementation should distinguish between the aligned
storage abstraction and the instruction set used to process that storage.

A future implementation could therefore provide multiple execution paths:

\begin{verbatim}
aligned buffer
|
+---- scalar implementation
|
+---- SSE implementation
|
+---- AVX implementation
|
+---- AVX2 implementation
\end{verbatim}

The appropriate implementation could then be selected according to
compile-time target information or runtime CPU feature detection.

This would make the alignment component more portable across different
processor architectures and configurations encountered in automotive
systems.


% ------------------------------------------------------------
\subsection{Generalisation Beyond \texttt{f32}}

The current \texttt{AlignedBuffer} is specifically designed around
\texttt{f32} values and the \texttt{F32x8} representation.

This design is appropriate for demonstrating AVX processing, but it means
that the buffer abstraction itself is tied to one data type.

A future implementation could separate the generic aligned storage
abstraction from the particular SIMD kernel. For example:

\begin{verbatim}
AlignedBuffer<T, A>
\end{verbatim}

could represent aligned storage independently of whether the stored values
are floating-point values, integers, or user-defined structures.

The SIMD operation could then be implemented separately for each supported
type and architecture.

Such a design would make the alignment component more closely resemble the
general-purpose nature of Boost.Align while retaining specialized optimized
implementations where appropriate.


% ------------------------------------------------------------
\subsection{Improved Memory and Performance Evaluation}

The current benchmark provides a positive performance result for the selected
AVX workload. The repository reports an average iteration time of
approximately 20.70~ms for Rust compared with 21.68~ms for Boost.Align.
The corresponding reported throughput is approximately 6.48~GB/s for Rust
and 6.19~GB/s for Boost.Align, giving a Boost-to-Rust execution-time ratio
of approximately 1.047. Thus, the Rust implementation was approximately
approximately 4.7\% lower measured execution time in the selected workload.

The result should nevertheless be interpreted as a workload-specific
observation rather than evidence that Rust is generally faster than
Boost.Align. The benchmark uses a single 64~MiB workload, a 32-byte
alignment requirement and AVX processing.

Furthermore, the two SIMD kernels are not instruction-identical. The C++
benchmark uses \texttt{\_mm256\_mul\_ps} with a vector containing
\texttt{2.0f}, whereas the Rust implementation uses
\texttt{\_mm256\_add\_ps(v,v)}. Although both operations produce the same
mathematical result for the benchmark input, they represent different
instructions and may have different performance characteristics.

The memory measurements also identify an important area for further
investigation. The reported RSS increase during allocation was approximately
196~KB for the Boost implementation and 65,540~KB for the Rust
implementation, while both implementations subsequently showed approximately
65,540~KB of RSS reduction during deallocation.

This difference should therefore be investigated before drawing conclusions
about allocator efficiency or memory footprint.

A more comprehensive evaluation should therefore vary:

\begin{itemize}
    \item buffer size;
    \item alignment;
    \item data type;
    \item CPU architecture;
    \item SIMD instruction set;
    \item compiler optimization level;
    \item allocation frequency; and
    \item number of repeated operations.
\end{itemize}

The benchmark should also use identical SIMD operations and equivalent
compiler settings for both implementations. In addition, the memory
measurements should be repeated independently from the processing benchmark
to determine whether the observed RSS difference is caused by allocation
behaviour, measurement methodology, allocator caching, or another property of
the experimental environment.

These changes would allow the performance conclusions to be generalized
beyond the specific AVX workload used in the current implementation.


% ============================================================
\section{Boost.Endian Improvements}

\label{sec:endian-improvements}

The Boost.Endian implementation provides the three principal usage models
identified in Boost.Endian: conversion functions, buffer types and arithmetic
types. The implementation also provides C++ interoperability through the
\texttt{cxx} crate.

The current implementation is therefore broader than the Boost.Align
implementation in terms of API structure. Nevertheless, several improvements
would increase its completeness, robustness and suitability for
interoperability.


% ------------------------------------------------------------
\subsection{Stabilisation of the Baseline Implementation}

The repository contains several Endian implementation variants developed
during the migration process. The comparison documentation identifies the
original \texttt{endian} crate as the Rust-native baseline, an
\texttt{endianclv1} implementation using a C ABI, an
\texttt{endianclv2} implementation using a \texttt{cxx} bridge, and an
\texttt{endianv3} implementation providing a richer demonstration interface.

This experimentation is valuable because it demonstrates that FFI design is
itself an important part of the migration process.

However, the baseline implementation should be consolidated before being
treated as the final migrated component. In particular, the repository
documentation records that the original Rust-native implementation still
requires further stabilization.

A future version should establish one canonical Endian implementation and
clearly separate experimental alternatives from the final migration result.

The final implementation should then be evaluated using a reproducible test
and benchmark procedure.


% ------------------------------------------------------------
\subsection{Expansion of Non-Standard Integer Widths}

Boost.Endian supports integer representations whose widths do not necessarily
correspond to standard C++ integer types. This includes 24-, 40-, 48- and
56-bit representations \cite{boostendian}.

The current Rust implementation provides explicit 24-bit load and store
functionality, with the value represented internally using a wider native
integer type.

A natural extension would be to generalize this mechanism to all non-standard
widths supported by Boost.Endian:

\begin{verbatim}
24-bit
40-bit
48-bit
56-bit
\end{verbatim}

The representation should continue to use byte-oriented storage rather than
relying on a native integer with a matching size, since Rust does not provide
native integer types for these widths.

This would allow the Rust implementation to cover a larger portion of the
Boost.Endian buffer API while retaining the portability advantages of
explicit byte representations.


% ------------------------------------------------------------
\subsection{Aligned and Unaligned Buffer Types}

The Boost.Endian buffer interface distinguishes aligned and unaligned
representations. The official Boost documentation identifies unaligned
buffers as useful for portable representations because they avoid padding
and architecture-dependent alignment requirements
\cite{boostendianbuffers}.

The current Rust implementation primarily represents endian-aware values
through transparent wrappers around primitive types.

A future implementation could make the distinction between aligned and
unaligned representations explicit.

An unaligned representation could use:

\begin{verbatim}
[u8; N]
\end{verbatim}

as its underlying storage.

An aligned representation could use an appropriately aligned wrapper where
the target architecture and use case justify the additional constraint.

This would allow the migration to reproduce another important design choice
in Boost.Endian while allowing applications to choose portability or
potentially more efficient aligned access according to their requirements.


% ------------------------------------------------------------
\subsection{Expansion of the C++ FFI Surface}

The current \texttt{cxx} bridge exposes a selected subset of the Rust
implementation. In particular, the demonstrated bridge includes load/store
and buffer functionality for 16-, 32- and 64-bit integer widths.

The Rust implementation also contains 24-bit load/store functionality that
is not currently exposed through the C++ bridge.

A future version should define an explicit FFI coverage matrix:

\begin{table}[H]
\centering
\renewcommand{\arraystretch}{1.2}
\begin{tabular}{|l|c|c|c|}
\hline
\textbf{Feature} & \textbf{Boost.Endian} &
\textbf{Rust API} & \textbf{C++ FFI} \\
\hline
16-bit conversion & Yes & Yes & Yes \\
24-bit conversion & Yes & Yes & No \\
32-bit conversion & Yes & Yes & Yes \\
40-bit representation & Yes & No & No \\
48-bit representation & Yes & No & No \\
56-bit representation & Yes & No & No \\
64-bit conversion & Yes & Yes & Yes \\
\hline
\end{tabular}
\caption{Current and potential Boost.Endian interoperability coverage}
\label{tab:endian-ffi-coverage}
\end{table}

This would make the interoperability scope explicit and prevent the C++
interface from being interpreted as a complete representation of the Rust
implementation.


% ------------------------------------------------------------
\subsection{Improved C++ Arithmetic Interface}

Rust implements the arithmetic functionality through Rust operator traits
such as \texttt{Add}, \texttt{Sub}, \texttt{Mul}, \texttt{Div},
\texttt{BitAnd}, \texttt{BitOr} and \texttt{BitXor}. These traits provide
an idiomatic Rust representation of Boost.Endian's arithmetic types.

However, operator traits cannot simply be exported through the FFI as C++
operator overloads.

The current solution therefore uses a C++ wrapper around the FFI functions.

A future implementation could improve this interface by providing a
dedicated C++ compatibility layer that exposes an API closer to the
corresponding Boost.Endian arithmetic types.

The C++ layer would remain thin, with the actual endian representation and
conversion logic remaining in Rust.

This would produce the following architecture:

\begin{verbatim}
C++ compatibility type
|
v
cxx bridge
|
v
Rust EndianInt
|
v
Encoding + conversion
\end{verbatim}

Such an arrangement would allow existing C++ applications to migrate
individual uses of Boost.Endian without requiring immediate changes to their
surrounding code.


% ------------------------------------------------------------
\subsection{Explicit Validation of Binary Layout}

Endian-aware data types are frequently used in binary protocols and file
formats. Therefore, numerical correctness alone is insufficient to validate
the implementation.

The tests should also verify the exact byte-level representation.

For example, a value:

\[
0x01020304
\]

stored as big endian must produce:

\begin{verbatim}
01 02 03 04
\end{verbatim}

while the corresponding little-endian representation must produce:

\begin{verbatim}
04 03 02 01
\end{verbatim}

The current C++ demonstration already exercises mixed-endian structures and
checks their binary representation. This should be expanded into a
systematic automated test suite.

The test suite should verify:

\begin{itemize}
    \item byte representation;
    \item structure size;
    \item field offsets;
    \item alignment;
    \item signed values;
    \item negative values for non-standard widths;
    \item floating-point representations;
    \item round-trip conversions; and
    \item mixed-endian structures.
\end{itemize}

These tests are particularly important because a numerically correct
conversion can still result in an incorrect binary layout.


% ------------------------------------------------------------
\subsection{Improved FFI Error Handling}

The current Endian FFI uses Rust slices to represent buffers crossing the
C++/Rust boundary. This is considerably safer than reconstructing arbitrary
Rust slices from raw pointers, but the caller is still responsible for
supplying a buffer with the expected number of bytes.

For example, a 64-bit load requires at least eight bytes.

A future FFI layer should therefore expose fallible versions of the load and
store operations. Instead of allowing an invalid slice length to trigger a
panic, the interface could return an explicit status:

\begin{verbatim}
SUCCESS
INVALID_LENGTH
INVALID_ARGUMENT
UNSUPPORTED_TYPE
\end{verbatim}

This would provide a more deterministic interface for C++ applications and
would be preferable for production-oriented integration.


% ------------------------------------------------------------
\subsection{FFI Boundary Optimisation}

The repository contains a dedicated benchmark comparing a native C++
implementation with the Rust implementation accessed through the
\texttt{cxx} bridge. The benchmark compares manual C++ big-endian 64-bit
load/store operations with the Rust \texttt{load\_big\_u64} and
\texttt{store\_big\_u64} functions.

Both implementations pass the correctness and byte-order checks. Across the
four tested input sizes, the Rust FFI path reports lower execution time and
higher throughput. The Rust times range from 880~ns to 60,191,448~ns per
iteration, while the C++ times range from 1,406~ns to 92,625,204~ns.

The reported throughput is approximately 4.16--4.66~GB/s for Rust compared
with 2.69--2.91~GB/s for C++, with an average reported speedup of
approximately 35\%.

The results are summarized in Table~\ref{tab:endian-benchmark-results}.

\begin{table}[H]
\centering
\renewcommand{\arraystretch}{1.2}
\begin{tabular}{|r|r|r|r|r|}
\hline
\textbf{Elements} &
\textbf{C++ ns/iter} &
\textbf{Rust ns/iter} &
\textbf{C++ GB/s} &
\textbf{Rust GB/s} \\
\hline
256 &
1,406 &
880 &
2.91 &
4.66 \\
\hline
16,384 &
97,331 &
62,995 &
2.69 &
4.16 \\
\hline
1,048,576 &
5,884,663 &
3,842,808 &
2.85 &
4.37 \\
\hline
16,777,216 &
92,625,204 &
60,191,448 &
2.90 &
4.46 \\
\hline
\end{tabular}
\caption{Boost.Endian FFI benchmark results}
\label{tab:endian-benchmark-results}
\end{table}

This result is encouraging because the Rust implementation remains
competitive even though the C++ caller crosses the \texttt{cxx} boundary.
However, the comparison is not a direct comparison with the Boost.Endian
implementation. The C++ reference is a manual byte-shift implementation,
while the Rust path calls the conversion functions through
\texttt{rust::Slice}. Consequently, the benchmark should be interpreted as
an evaluation of the selected Rust FFI implementation against a native C++
byte-manipulation baseline, rather than as proof that Rust is faster than
Boost.Endian.

The benchmark should therefore distinguish between two different sources of
overhead:

\begin{enumerate}
    \item the cost of the endian conversion itself; and
    \item the cost of invoking the Rust implementation through FFI.
\end{enumerate}

The current benchmark performs a Rust load and Rust store for each element.
Consequently, the number of FFI calls grows with the number of processed
values.

Although this path performs well in the tested environment, a more
representative systems-level implementation should batch operations so that
one FFI call processes an entire buffer.

For example:

\begin{verbatim}
C++ buffer
|
| one FFI call
v
Rust process_buffer()
|
| process N elements
v
C++ buffer
\end{verbatim}

This would move the FFI boundary outside the inner processing loop and make
the cost of the interface easier to separate from the cost of the endian
conversion itself.

A future evaluation should also include a direct Boost.Endian baseline and
report the results separately for conversion cost, FFI cost, and complete
application-level processing.


% ============================================================
\section{Common FFI Improvements}

\label{sec:ffi-improvements}

Although Boost.Align and Boost.Endian have different implementation
characteristics, both expose similar challenges at the C++/Rust boundary.

The first improvement is to define the ownership model explicitly for every
FFI object. Any pointer crossing the boundary should have a documented
owner, lifetime and destruction mechanism.

The second improvement is to avoid allowing Rust panics to become the normal
error-reporting mechanism for invalid C++ input. FFI functions should instead
return explicit status values or use a documented error representation.

The third improvement is to minimize the amount of data exchanged through the
FFI boundary. Large buffers should preferably remain owned by the caller
while Rust operates on borrowed slices or pointers whose validity is
established by the interface contract.

The fourth improvement is to keep unsafe code localized inside the FFI layer.
Safe Rust abstractions should validate the relevant invariants before
invoking the low-level operations.

This produces the general migration boundary:

\begin{verbatim}
+----------------------+       +----------------------+
|      C++ system      |       |     Rust library     |
|                      |       |                      |
| Existing application | <---> | Safe abstractions    |
| Existing data types  |  FFI  | Type-level invariants|
|                      |       | Low-level operations|
+----------------------+       +----------------------+
             \                   /
              \                 /
               +---------------+
                 FFI boundary
             ownership/errors/
              representation
\end{verbatim}

Such a boundary is consistent with the migration objectives defined earlier
in this dissertation, particularly the requirement to maintain C++
interoperability while minimizing and isolating unsafe Rust.


% ============================================================
\section{Improvements to Validation}

\label{sec:future-validation}

The current validation demonstrates functional correctness, binary
representation, alignment and selected performance characteristics.

The benchmark results provide initial quantitative evidence for the
performance objectives of the migration strategy. The Boost.Align benchmark
measured approximately 20.70~ms per iteration for Rust versus 21.68~ms for
Boost.Align, corresponding to approximately 4.7\% lower measured execution
time for the Rust implementation. The Endian FFI benchmark reported an
average speedup of approximately 35\% for Rust against its manual C++
baseline.

These results are useful but remain specific to the tested workloads and
environment. A more comprehensive evaluation should therefore extend the
validation in several dimensions.


% ------------------------------------------------------------
\subsection{Cross-Architecture Testing}

Both selected libraries interact closely with processor architecture.

For Boost.Align, this concerns alignment requirements and SIMD instruction
availability.

For Boost.Endian, this concerns native byte order and the relationship
between the external representation and the host representation.

The implementation should therefore be evaluated on both little-endian and
big-endian architectures where practical.

This would verify that the implementation does not accidentally depend on
the little-endian architecture used by the development environment.


% ------------------------------------------------------------
\subsection{Compiler and Toolchain Testing}

The C++ interoperability layer should also be evaluated using multiple
compiler and toolchain configurations.

At minimum, the evaluation should consider:

\begin{itemize}
    \item GCC;
    \item Clang;
    \item different supported Rust toolchains;
    \item debug and release builds; and
    \item different optimization levels.
\end{itemize}

This is particularly relevant because Boost libraries were historically
designed to compensate for differences between C++ compilers and standard
library implementations.

A Rust migration should therefore demonstrate that the new implementation
does not simply replace one platform-specific dependency with another.


% ------------------------------------------------------------
\subsection{Property-Based Testing}

The existing tests use specific representative values. These are useful for
validating known cases but do not exhaustively explore the possible input
space.

Property-based testing could provide stronger validation.

For endian conversion, an important property is:

\[
reverse(reverse(x)) = x
\]

for all supported values \(x\).

Similarly, for a supported endian representation \(E\):

\[
E^{-1}(E(x)) = x
\]

should hold for all valid values.

For aligned buffers, properties can include:

\[
address \bmod alignment = 0
\]

and:

\[
length \bmod SIMD\_width = 0
\]

for buffers passed to the SIMD implementation.

These properties can complement the existing example-based tests.


% ------------------------------------------------------------
\subsection{Benchmark Reproducibility}

The benchmark results demonstrate the importance of reproducibility in
performance comparisons.

Future measurements should record the complete experimental configuration,
including:

\begin{itemize}
    \item processor model;
    \item number of available CPU cores;
    \item operating system;
    \item compiler versions;
    \item Rust toolchain version;
    \item optimization flags;
    \item SIMD features;
    \item buffer size;
    \item number of iterations; and
    \item allocation strategy.
\end{itemize}

Performance measurements should ideally be repeated sufficiently many times
to report not only the average but also measures such as standard deviation,
minimum and maximum execution time.

This would provide stronger evidence regarding whether the observed
differences are stable and reproducible.

For Boost.Align, the C++ and Rust kernels should additionally use equivalent
instructions. The current benchmark uses multiplication by \texttt{2.0f}
in C++ and vector addition in Rust. A future benchmark should use the same
operation on both sides so that the comparison isolates the implementation
and language differences more effectively.

For Boost.Endian, the benchmark should additionally include a direct
Boost.Endian implementation. This would allow three separate comparisons:

\begin{enumerate}
    \item native C++ byte manipulation;
    \item Boost.Endian; and
    \item Rust Endian through FFI.
\end{enumerate}

Such a comparison would provide stronger evidence for the performance
implications of the proposed migration.


% ============================================================
\section{Future Integration with Automotive Software}

The selected improvements also provide a path toward evaluating the migrated
libraries in a more realistic automotive software environment.

For Boost.Align, future work could investigate aligned buffers in workloads
involving:

\begin{itemize}
    \item sensor processing;
    \item SIMD-based signal processing;
    \item DMA-compatible memory;
    \item memory pools; and
    \item performance-critical data structures.
\end{itemize}

For Boost.Endian, representative workloads could include:

\begin{itemize}
    \item network protocol parsing;
    \item diagnostic messages;
    \item serialization;
    \item binary file formats; and
    \item communication between heterogeneous processors.
\end{itemize}

These scenarios would allow the migration methodology to be evaluated beyond
isolated library benchmarks and closer to the types of low-level
functionality encountered in automotive software architectures.

The Eclipse S-CORE context described earlier in this dissertation provides
one potential environment in which these evaluations could eventually be
performed.

The benchmark results also indicate that future automotive evaluation should
consider not only raw execution time but also the location of the
interoperability boundary. For high-frequency operations, reducing the number
of FFI calls may be more important than optimizing an individual conversion
operation. Similarly, for SIMD-oriented workloads, the interaction between
alignment guarantees, instruction-set availability and memory allocation
should be evaluated as a combined system property.


% ============================================================
\section{Future Work Summary}

The implementations presented in Chapter~\ref{chap:Chapter6} demonstrate the feasibility of
reimplementing selected Boost functionality in Rust, and the benchmark
results provide initial quantitative evidence for the performance objectives
of the migration strategy.

For Boost.Align, the benchmark measured approximately 20.70~ms per iteration
for Rust versus 21.68~ms for Boost.Align. The corresponding throughput was
approximately 6.48~GB/s for Rust and 6.19~GB/s for Boost.Align, resulting in
approximately 4.7\% lower measured execution time for the Rust
implementation. However, the different SIMD instructions and the observed
difference in RSS behaviour mean that the result should be considered
workload-specific rather than evidence of general Rust performance
superiority.

For Boost.Endian, the FFI benchmark reported lower execution times for Rust
at all tested input sizes, with reported throughput between approximately
4.16 and 4.66~GB/s compared with 2.69--2.91~GB/s for the manual C++ baseline.
The reported average speedup was approximately 35\%. This result demonstrates
that the selected Rust implementation can provide competitive performance
despite the use of a C++/Rust interoperability layer. However, because the
reference implementation is a manual C++ byte-shift implementation rather
than Boost.Endian itself, the result does not establish that Rust is faster
than Boost.Endian.

For Boost.Align, the principal improvements are the generalisation of
alignment requirements, expansion of the API, reduction of unsafe slice
construction, improved FFI error handling, runtime SIMD feature handling,
support for additional data types, and broader performance evaluation.

For Boost.Endian, the principal improvements are stabilisation of the
selected implementation variant, expansion of non-standard integer widths,
explicit aligned and unaligned representations, broader C++ FFI coverage,
improved arithmetic interoperability, stronger binary-layout validation,
explicit FFI error handling, and reduction of per-element FFI overhead.

Across both libraries, the benchmark results suggest that future work should
focus on improving the experimental methodology as well as the
implementations themselves. In particular, identical workloads and
instructions should be used when comparing implementations, direct Boost
baselines should be added where they are currently absent, and the cost of
repeated FFI crossings should be separated from the cost of the underlying
operations.

The most important architectural improvement remains the definition of a
robust interoperability boundary. Ownership, representation, lifetime,
error handling and safety invariants should be explicit at this boundary.

These improvements provide a path from the current research implementation
toward a more general and production-oriented Rust library while preserving
the central objective of the dissertation: demonstrating how low-level Boost
functionality can be migrated incrementally to Rust without requiring an
immediate replacement of the surrounding C++ system.